%!TEX root = ../nauchka.tex

\section{Выбор метода параллелизма}

\begin{frame}
\frametitle{Сравнение подходов}

\begin{table}[h]
\centering
\scriptsize
\begin{tabular}{|l|c|c|c|}
\hline
\textbf{Характеристика} & \textbf{TP} & \textbf{PP} & \textbf{FSDP} \\
\hline
Память на GPU & $O(M/P)$ & $O(M/P)$ & $O(M)$ \\
\hline
Коммуникации & Высокие & Высокие & Низкие \\
\hline
Шардирует $A$-матрицы & Да & Нет & Нет \\
\hline
Подходит для BaB & Да & Нет & Нет \\
\hline
Сложность реализации & Высокая & Очень высокая & Средняя \\
\hline
\end{tabular}
\end{table}

\vspace{0.5cm}

\textbf{Вывод:} Tensor Parallelism --- оптимальный выбор для верификации

\end{frame}

\begin{frame}
\frametitle{Tensor Parallelism: Column Parallel}

\begin{block}{Column Parallel слой}
Веса разделены по \textbf{выходной} размерности
\end{block}

\vspace{0.3cm}

\textbf{Прямой проход:} $Y = X \cdot W$
\begin{itemize}
    \item $W = [W_0 | W_1]$ разделен между GPU
    \item $Y = [Y_0 | Y_1]$ также разделен
    \item Требуется AllGather для получения полного $Y$
\end{itemize}

\vspace{0.3cm}

\textbf{Обратный CROWN:} $A_{prev} = A_{curr} \cdot W$
\begin{itemize}
    \item $A_{curr}$ уже разделена: $[A_0 | A_1]$
    \item Локально: $P_i = A_i \cdot W_i$ (частичное произведение)
    \item \textbf{AllReduce:} $A_{prev} = P_0 + P_1$ на всех GPU
\end{itemize}

\end{frame}

\begin{frame}
\frametitle{Tensor Parallelism: Row Parallel}

\begin{block}{Row Parallel слой}
Веса разделены по \textbf{входной} размерности
\end{block}

\vspace{0.3cm}

\textbf{Прямой проход:} $Y = X \cdot W$
\begin{itemize}
    \item $X = [X_0 | X_1]$ разделен между GPU
    \item Локально: $P_i = X_i \cdot W_i$
    \item AllReduce: $Y = P_0 + P_1$
\end{itemize}

\vspace{0.3cm}

\textbf{Обратный CROWN:} $A_{prev} = A_{curr} \cdot W$
\begin{itemize}
    \item $A_{curr}$ реплицирована (полная матрица на всех GPU)
    \item Локально: $A_{prev,i} = A_{curr} \cdot W_i$
    \item Результат автоматически разделен
    \item \textbf{Коммуникация не нужна!}
\end{itemize}

\end{frame}

\begin{frame}
\frametitle{Почему не Pipeline Parallelism?}

\begin{alertblock}{Проблемы PP для верификации}
\begin{enumerate}
    \item \textbf{Skip connections} в ResNet/Transformer требуют передачи матриц $A$ между GPU
    \item \textbf{Двунаправленные зависимости:}
    \begin{itemize}
        \item Прямой проход: для вычисления границ ReLU
        \item Обратный CROWN: для распространения $A$
    \end{itemize}
    \item \textbf{Bubbles:} PP эффективен только при большом количестве микро-батчей, но в верификации одного изображения такого нет
    \item НЕ шардирует матрицы $A$ --- проблема памяти остается!
\end{enumerate}
\end{alertblock}

\vspace{0.3cm}

\textbf{Вывод:} PP применим только для очень глубоких сетей без skip connections

\end{frame}

\begin{frame}
\frametitle{Почему не FSDP?}

\begin{alertblock}{Проблемы FSDP для верификации}
\begin{enumerate}
    \item В режиме inference (верификация) нет:
    \begin{itemize}
        \item Градиентов
        \item Состояний оптимизатора
    \end{itemize}
    \item Веса модели обычно малы по сравнению с матрицами $A$
    \item FSDP НЕ шардирует активации и промежуточные матрицы $A$
    \item Решает проблему хранения весов, но игнорирует главную проблему верификации
\end{enumerate}
\end{alertblock}

\vspace{0.3cm}

\textbf{Вывод:} FSDP не решает проблему памяти для матриц $A$

\end{frame}
